\section{The discriminant form prime by prime}

Throughout this section $E_p=\coker(G)_p$ with the pairing $\beta_G$, and
$\ell_p$ denotes length. The maximal totally isotropic subgroup of $E_p$ bounds
$\Delt_p$ from above; the square-order bound of Proposition~\ref{prop:gp} is the
statement that this maximum is at most $|E_p|^{1/2}$.

\subsection{Good primes}

\begin{proposition}[good-prime decomposition]\label{prop:dp}
For every prime $p\nmid n$, the pointed lattice splits $\beta_G$-orthogonally as
\[
E_p\ \cong\ \big(W_p/\alpha W_p,\ (x,y)\mapsto x\!\cdot\! y/\alpha\big)\ \perp\
\big(\coker K_p,\ K_p^{-1}\big),\quad
K_p=\begin{pmatrix}n(rk-\lambda)&nr\\ nr&b\end{pmatrix},
\]
where $W_p=\{x\in\ZZ_p^{n}:\mathbf 1^{T}x=0\}$; here $\det K_p=n\alpha r/k$.
\end{proposition}

\begin{proof}
Since $p\nmid n$, the vector $\mathbf 1_n$ has unit norm $n$, so
$\ZZ_p^{n}=W_p\perp\ZZ_p\mathbf 1_n$; adjoining the apex coordinate $\ZZ_p e_*$
gives the orthogonal splitting $\ZZ_p^{v}=W_p\perp T_p$ with
$T_p=\ZZ_p\mathbf 1_n\oplus\ZZ_p e_*$. Restricting $G$ from
Proposition~\ref{prop:gp}: on $W_p$ (coordinate sum zero) it is $\alpha$ times
the dot product, which is unimodular there; on $T_p$ it is the block $K_p$, whose
determinant $n\alpha r/k$ follows from $(rk-\lambda)b-nr^{2}=\alpha r/k$. Passing
to cokernels and inverse-Gram pairings in this orthogonal basis gives the stated
sum; $\ell_p$ of each summand is read from the entry valuations of $K_p$.
\end{proof}

\subsection{The odd-prime Witt bound}

The determinant class distinguishes the two possible Witt indices in even
dimension. The following calculation makes that class explicit.

\begin{theorem}[odd valuation-one Witt obstruction]\label{thm:wi}
Let $p$ be odd with $p\nmid k\lambda$ and $v_p(\alpha)=1$, and suppose $n=v-1$ is
even, $n=2m$. Then $E_p\cong\FF_p^{\,n}$ carries a nondegenerate symmetric form,
its dual selected defect is totally isotropic, and its minimal number of
generators $\mu(\Delt_p)$ satisfies
\[
\mu(\Delt_p)\le m\quad\text{if }(-1)^m k\text{ is a square mod }p,\qquad
\mu(\Delt_p)\le m-1\quad\text{otherwise.}
\]
The second line is one generator stronger than $|\Delt_p|^{2}\mid\det G$. For
$v=5+60t$ with $t\ge1$ and $3\nmid t$, Hanani's theorem~\cite{Hanani1972}
supplies an $S(2,5,v)$; then $n=4+60t$ is even, at $p=3$ one has $m=2+30t$,
$(-1)^m k=5$ a nonsquare, so $\Delt_3\cong(\ZZ/3)^{\mu}$ with $\mu\le 1+30t$,
whereas the coarse envelope permits $2+30t$. This is an infinite strict
obstruction.
\end{theorem}

\begin{proof}
The hypotheses give $r\equiv\lambda\pmod p$. The identities
$r(k-1)=\lambda(v-1)$ and $bk=vr$ then give $v\equiv k$ and
$b\equiv\lambda\pmod p$, respectively. Thus $G\equiv\lambda J_v\pmod p$ has rank one, and
$v_p(\det G)=n$ forces every nonunit invariant factor to have valuation one, so
$E_p\cong\FF_p^{n}$. Identify $(1/p)\ZZ/\ZZ$ with $\FF_p$ by
$a/p\mapsto a\bmod p$. Direct inversion gives
\[
G^{-1}=\frac1\alpha
\begin{pmatrix}I_n+J_n&-k\mathbf1_n\\
-k\mathbf1_n^T&k(kr-\lambda)/r\end{pmatrix}.
\]
The first $n$ coordinate classes form a basis of $E_p$: the image of $G$
modulo $p$ is the constant line in $\FF_p^v$. In this basis the induced
finite-field form has matrix $u^{-1}(I_n+J_n)$, where $u=\alpha/p$.
Its determinant is $v/u^n\equiv k/u^n\pmod p$, of square class $k$ for even $n$.

By the classical determinant criterion~\cite[Theorem 5.2]{WilsonNotes}
(see also~\cite{Schmidt}), a
$2m$-dimensional such space has Witt index $m$ iff $(-1)^m$ times its determinant
is a square, else $m-1$; the dual defect is a totally isotropic subspace, so its
dimension is at most the Witt index. The $S(2,5,v)$ evaluation is direct.
\end{proof}

\subsection{Exceptional primes}

For $p\mid n$ the root/line split fails; a different primitive orthogonal
splitting survives.

\begin{theorem}[exceptional-prime reduction]\label{thm:xd}
For every prime $p\mid n$ the pointed lattice is $\ZZ_p$-congruent to
$\alpha D_{n-2}\perp C_p$, where $D_{n-2}=I+J$ has rank $n-2$ and determinant
$n-1$ (a $p$-adic unit, so $D_{n-2}$ is unimodular over $\ZZ_p$), and
\[
C_p=\begin{pmatrix}\alpha n/(n-1)&-\alpha/(n-1)&0\\
-\alpha/(n-1)&\lambda+\alpha/(n-1)&r\\ 0&r&b\end{pmatrix},\qquad
\det C_p=\frac{\alpha^{2}r}{k(n-1)} .
\]
Hence all but at most three local invariant factors have valuation
$v_p(\alpha)$, and the remaining three are read from the determinantal ideals of
$C_p$. At $p=2$ the root block $\alpha D_{n-2}$ is type II; if $v_2(\alpha)=1$
and $k\lambda$ is odd, the whole discriminant space is symplectic of dimension
$n$, so bilinear isotropy meets---but does not beat---the square-order bound.
\end{theorem}

\begin{proof}
Put $t=n-1$, a $p$-adic unit. The sublattice
$U=\{x\in\ZZ_p^{n}:x_n=0,\ \sum_i x_i=0\}$ has dot-product Gram $D_{n-2}=I+J$
($\det=n-1$, unimodular over $\ZZ_p$) and is $G$-orthogonal to
$s=(\mathbf 1_n-e_n)/(n-1)$, $a=e_n-s$, and the apex $e_*$. A direct computation
gives $\langle a,a\rangle_G=\alpha n/(n-1)$, $\langle a,s\rangle_G=-\alpha/(n-1)$,
$\langle s,s\rangle_G=\lambda+\alpha/(n-1)$, $\langle a,e_*\rangle_G=0$, and
$\langle s,e_*\rangle_G=r$, so in the local basis
$(a,s,e_*)$ the residual Gram is exactly $C_p$; the determinant identity
$\alpha^{n-2}\det D_{n-2}\det C_p=\alpha^{n}r/k=\det G$ together with adjugation
gives $\det C_p$ and $C_p^{-1}$, and the three residual exponents are the
successive determinantal-ideal valuations. This is a unimodular local change of
basis, not the invalid zero-sum-root/constant-line split (whose index is
divisible by $p$).

At $p=2$, $D_{n-2}=I+J$ has even diagonal and odd determinant,
hence gives a type-II $2^{v_2(\alpha)}$-modular constituent (standard $\ZZ_2$
Jordan splitting~\cite{ConwaySloane}); in the valuation-one unit case
$2G^{-1}\bmod2$ is alternating, so the discriminant space is symplectic and its
maximal isotropic dimension equals the square-order bound.
\end{proof}

\subsection{The binary order bound}

For subgroup orders the binary case has a simpler general answer, independent
of the design. Write $\mathrm{maxiso}(E)$ for the maximum order of a totally
isotropic subgroup of a finite bilinear group $E$. A nonsingular group is
\emph{metabolic} if it has an isotropic subgroup of order $|E|^{1/2}$; this
does not require that subgroup to be a direct summand.

\begin{lemma}[binary bilinear order bound]\label{lem:plateau}
Let $E$ be a finite $2$-group with a nonsingular symmetric bilinear pairing
$\beta:E\times E\to\QQ/\ZZ$, and put $\ell=\log_2|E|$. Then
\[
\mathrm{maxiso}(E)=2^{\lfloor\ell/2\rfloor}.
\]
In particular every such group of square order is metabolic.
\end{lemma}

\begin{proof}
For every isotropic subgroup $S$, nonsingularity gives
$|S|\,|S^\perp|=|E|$ and $S\le S^\perp$, proving the upper bound.

If $|E|\ge4$, there is an isotropic involution. Indeed, if
$\dim_{\FF_2}E[2]\ge2$, the map $x\mapsto\beta(x,x)$ from $E[2]$ to
$(1/2)\ZZ/\ZZ$ is linear, since $2\beta(x,y)=0$, so it has a nonzero kernel.
Otherwise $E$ is cyclic of order $2^s$, $s\ge2$; for a generator $g$ with
$\beta(g,g)=a/2^s$, its involution has self-pairing
$\beta(2^{s-1}g,2^{s-1}g)=a2^{s-2}=0$ in $\QQ/\ZZ$.

Let $H$ be the isotropic subgroup of order two so obtained.
The induced pairing on $H^\perp/H$ is nonsingular and this quotient has order
$|E|/4$. Induction, starting at orders one and two, gives an isotropic subgroup
there of order $2^{\lfloor(\ell-2)/2\rfloor}$. Its inverse image in $H^\perp$
is isotropic of order $2^{\lfloor\ell/2\rfloor}$.
\end{proof}

\begin{corollary}[binary order consequence]\label{thm:xi}
For every design under consideration,
\[
\mathrm{maxiso}(E_2)=
2^{\left\lfloor(nv_2(\alpha)+v_2(r)-v_2(k))/2\right\rfloor}.
\]
Thus bilinear isotropy gives no smaller bound on the \emph{order} of
$\Delt_2$ than the rounded square-divisibility bound. In particular, if
$2\mid n$ and $A=v_2(\alpha)\ge1$, the decomposition of
Theorem~\ref{thm:xd} satisfies
\[
\mathrm{maxiso}(E_2)=2^{A(n-2)/2}\,
\mathrm{maxiso}(\coker C_2).
\]
\end{corollary}

\begin{proof}
Apply Lemma~\ref{lem:plateau} to $E_2$ and use
$\ell_2(E_2)=v_2(\det G)$. In the exceptional decomposition the root
constituent has even length $A(n-2)$, so the formula for its orthogonal sum
with the residual group gives the second equality.
\end{proof}

The rank-three reduction remains useful for invariant factors and the pairing,
but is not needed for this order conclusion. Lemma~\ref{lem:plateau} is an
elementary fact about finite bilinear groups, proved here for clarity; it is
not a design-specific obstruction. It does not classify isotropically
embeddable subgroup types or quadratic refinements.

Metabolicity also must
not be confused with hyperbolicity: the plane
$V(4)=\frac14\left[\begin{smallmatrix}2&1\\1&2\end{smallmatrix}\right]$
on $(\ZZ/4)^2$ has the Lagrangian $2(\ZZ/4)^2$, but has no isotropic element
of order four and hence no free Lagrangian.
